% Adaptive, redundancy-aware probe construction
\begin{figure}[t]
\centering
\begin{tikzpicture}[crit/.style={circle, draw=lexurslate, fill=lexurslate!8,
   minimum size=8mm, font=\small, inner sep=0pt}]
% ---------- left: correlation graph + clusters ----------
\node[crit] (f1) at (0,0) {$f_1$};
\node[crit] (f2) at (1.6,0.9) {$f_2$};
\node[crit] (f3) at (2.5,-0.2) {$f_3$};
\node[crit] (f4) at (1.1,-1.5) {$f_4$};
\node[crit] (f5) at (2.9,-1.7) {$f_5$};
% redundant (positively correlated) edges, weight >= theta
\draw[lexurblue, line width=1.1pt] (f2) -- (f3);
\draw[lexurblue, line width=1.1pt] (f4) -- (f5);
\node[lexurblue, font=\scriptsize] at (2.15,0.55) {$\Corr\!\ge\!\theta$};
% cluster ellipses
\begin{scope}[on background layer]
\draw[dashed, lexurblue, rounded corners=10pt]
   ($(f2)+(-0.6,0.6)$) rectangle ($(f3)+(0.6,-0.6)$);
\draw[dashed, lexurblue, rounded corners=10pt]
   ($(f4)+(-0.6,0.6)$) rectangle ($(f5)+(0.6,-0.6)$);
\end{scope}
\node[font=\scriptsize\bfseries, lexurslate] at (1.4,-2.7) {correlation graph};
\node[lexurblue, font=\scriptsize] at (1.9,1.75) {clusters $C_1,C_2$};

% ---------- arrow ----------
\node[single arrow, draw=lexurslate, fill=lexurslate!10, minimum height=10mm,
      inner sep=2pt, font=\scriptsize] at (4.5,-0.4) {build};

% ---------- right: resulting probe family ----------
\begin{scope}[xshift=5.9cm, yshift=0.9cm,
      pr/.style={rounded corners=2pt, inner sep=3pt, font=\scriptsize, align=left}]
\node[pr, draw=lexurslate, fill=lexurslate!6, anchor=west] (s1) at (0,0)
   {$m$ singletons $q_i=r_i$ \scriptsize(separation)};
\node[pr, draw=lexurblue, fill=lexurblue!8, anchor=west] (s2) at (0,-0.75)
   {per cluster: $q_{C_k,\mathrm{mean}},\,q_{C_k,\max}$};
\node[pr, draw=lexurslate, fill=lexurslate!6, anchor=west] (s3) at (0,-1.5)
   {grand mean $q_\mu$, grand max $q_{\max}$};
\node[pr, draw=lexurcrimson, fill=lexurcrimson!8, anchor=west] (s4) at (0,-2.3)
   {$|\Q|=O(m+c)$ \scriptsize— linear, not exponential};
\node[font=\scriptsize\bfseries, lexurslate, above=1mm of s1] {declared probe family $\Q$};
\end{scope}
\end{tikzpicture}
\caption{Adaptive probe construction. Positively correlated (redundant) criteria
are grouped into clusters by thresholding the signed correlation at $\theta$.
Each nontrivial cluster contributes one mean and one max aggregate in place of
the exponential subcoalition aggregates; singleton probes remain for separation.
Together with the two grand probes, the family has $m+2+2c$ probes, where $c$ is
the number of nontrivial clusters.}
\label{fig:probes}
\end{figure}
